\SourceLeaf{114}
\pdfbookmark[1]{Göttingische gelehrte Anzeigen. 23 December 1815.}{anzeige1815}
\begin{center}
\rule{0.65\linewidth}{0.4pt}\par
Göttingische gelehrte Anzeigen. 23 December 1815.\par
\rule{0.65\linewidth}{0.4pt}
\end{center}

On 7 December, Professor \textsc{Gauss} submitted a lecture to the Royal Society, entitled:

\begin{center}
\emph{Another new proof of the theorem that every polynomial in one variable can be resolved into real factors of the first or second degree.}
\end{center}

The theorem stated here, the most important in the theory of equations, has, as is well known, received extensive attention from several of the foremost mathematicians. The principal attempts to give a fully rigorous proof of it, by \textsc{d'Alembert}, \textsc{Euler}, \textsc{Foncenex} and \textsc{Lagrange}, were already brought together and examined by the author of the present paper in a work published sixteen years ago. However much the ingenuity that distinguishes those works deserves recognition, it cannot be denied that they all leave more or fewer gaps that destroy their force as proofs. Although \textsc{Lagrange}'s penetrating investigations have remedied most of those defects, neither the efforts of this great mathematician nor the ingenious way in which \textsc{Laplace} later treated the subject can be exempted from the principal criticism that applies to all those attempted proofs. This is that the matter was treated as though only the \emph{form} of the roots had to be determined, while their \emph{existence} was assumed without proof,
\SourceLeaf{115}
a way of reasoning that is illusory in precisely this case and is in fact a true \emph{petitio principii}. All the proofs mentioned are purely analytic, including that of \textsc{d'Alembert}, although it appears in geometric clothing that is entirely irrelevant to the matter itself. One might therefore almost be tempted to conclude that this assumption, whose inadmissibility incidentally escaped the notice of that analyst himself, could not have been avoided in an analytic treatment.

In the work mentioned above, the author of the present paper had treated the subject in an entirely different way and given a very simple new proof, whose distinctive feature is that it rests in part on the geometry of position and otherwise seems to leave nothing to be desired in respect of rigour. At the same time, however, he had already declared that he by no means considered it impossible to arrive at a fully rigorous proof by purely analytic means, and had reserved a detailed development of his ideas for another time. Other occupations had until now drawn him away from this subject: the present paper is intended to fulfil that promise.

As for this new proof itself, it is certainly far inferior to the earlier one in simplicity and brevity; but this circumstance lies in the nature of the subtle investigation itself. It is always gratifying, alongside a very simple proof of the important theorem that rests in part on material from another field, to have a second proof that is admittedly longer and more elaborate, but entirely independent in itself. An account of the main ideas of the new proof, if it were to be even moderately satisfactory, would take up far more space than the purpose of these pages allows. We must therefore refer those who find a principal charm of mathematics in the perfect clarity of its truths and the utmost rigour of its proofs to the third volume of the Society's Commentationes, which will appear shortly.
\begin{center}\rule{0.28\linewidth}{0.4pt}\end{center}
